\documentclass[10pt,a4paper]{amsart}
\usepackage{amsmath,amssymb,amsthm,graphicx}
\usepackage{amsmath,amssymb,amsthm,xfrac,xspace,hyperref,graphicx,comment}
\excludecomment{DefTralics}
\providecommand{\C}{}
\providecommand{\bibliofont}{\small}
\newcommand{\ie}{i.e.,\xspace}
\newcommand{\eg}{e.g.,\xspace}
\newcommand{\etc}{etc.\xspace}
\newcommand{\cf}{cf.\xspace}
\newcommand{\resp}{resp.\xspace}
\providecommand{\smaller}{\small}
\newcommand{\Subsection}{\subsection}
\newcommand{\Subsubsection}{\subsubsection}
\newcommand{\skpt}{\smallskip}
\emergencystretch=3em
\setcounter{section}{1}
\providecommand{\noopsort}[1]{}

\multlinegap0pt
\usepackage{mathrsfs}
\let\mathcal\mathscr
\newenvironment{enumeratea}
{\bgroup\def\theenumi{\alph{enumi}}\begin{enumerate}}
{\end{enumerate}\egroup}

\newcommand\mto{\mathchoice{\longmapsto}{\mapsto}{\mapsto}{\mapsto}}
\newcommand\hto{\mathrel{\lhook\joinrel\to}}
\newcommand\To[1]{\mathchoice{\xrightarrow{\textstyle\kern4pt#1\kern3pt}}{\stackrel{#1}{\longrightarrow}}{}{}}
\newcommand{\sbullet}{{\scriptscriptstyle\bullet}}
\newcommand{\tprod}{{\textstyle\prod}}
\newcommand{\rc}{\mathrm{c}}


\usepackage[all,cmtip]{xy}\let\labelstyle\textstyle

\theoremstyle{plain}
\newtheorem{theo}{Theorem}[section]
\newtheorem*{theo*}{Theorem}
\newtheorem{prop}[theo]{Proposition}
\newtheorem{lem}[theo]{Lemma}
\newtheorem{cor}[theo]{Corollary}

\theoremstyle{definition}
\newtheorem{defi}[theo]{Definition}
\newtheorem{rema}[theo]{Remark}
\newtheorem{conj}[theo]{Conjecture}

\DeclareFontFamily{U}{wncy}{}
\DeclareFontShape{U}{wncy}{m}{n}{%
<5>wncyr5%
<6>wncyr6%
<7>wncyr7%
<8>wncyr8%
<9>wncyr9%
<10>wncyr10%
<11>wncyr10%
<12>wncyr6%
<14>wncyr7%
<17>wncyr8%
<20>wncyr10%
<25>wncyr10}{}
\DeclareMathAlphabet{\cyrille}{U}{wncy}{m}{n}

\DeclareMathOperator \br {Br}
\DeclareMathOperator \cok {coker}
\DeclareMathOperator\Cone{Cone}
\DeclareMathOperator \Imm {Im}\let\Im\Imm
\DeclareMathOperator \Ker {Ker}
\DeclareMathOperator \pic {Pic}
\DeclareMathOperator \gal {Gal}
\DeclareMathOperator \spec {Spec}

\renewcommand\AA{\mathbf{A}}
\newcommand \ZZ {\mathbf{Z}}
\newcommand \QQ {\mathbf{Q}}
\newcommand \FF {\mathbf{F}}
\newcommand \LL {\mathbf{L}}
\newcommand \NN {\mathbf{N}}
\newcommand \RR {\mathbf{R}}
\newcommand \Gm {\mathbf{G}_m}\let\G\Gm

\newcommand\cF {\mathcal{F}}
\newcommand\cN {\mathcal{N}}
\newcommand\cO {\mathcal{O}}
\newcommand \cT {\mathcal{T}}
\newcommand \cC {\mathcal{C}}

\newcommand\Sha{\cyrille X}

\let\ov\overline

\newcommand \kbar {{\bar k}}
\newcommand \et {\textup{ét}}
\newcommand \nr {\mathrm{nr}}
\newcommand \tors {\mathrm{tors}}
\title{Finiteness of the cohomology of tori and of finite multiplicative-type group schemes over a curve}
\author{Cyril Demarche, David Harari}
\date{doi:10.5802/jep.129; CC BY4.0}
\begin{document}
\maketitle
\noindent\textit{Editorial note.} All delivered mathematical text is the author's own, copied from the pinned publisher source sources/jep-2020-129-tori-duality/full-original-native.tex; no statement or proof step was rewritten, reordered or removed.\par\medskip
\noindent\textit{Editorial note.} Preamble provenance: the retained preamble is the source's own layout adapter plus the author's macro block at lines 10-88 of sources/jep-2020-129-tori-duality/full-original-native.tex (everything between the publisher class metadata and \textbackslash{}begin\{document\}, comments excluded). Nothing was added, no macro was invented, and the publisher document class is deliberately not reproduced; the excerpt serves only as preamble, never as a span coordinate.\par\medskip
\noindent\textit{Editorial note.} External references are navigation pointers only; every deliberately kept step of the author's argument stays inside the counted spans. The displayed numbers of those pointers are the author's own numbering under the author's \textbackslash{}newtheorem declarations (the source has no subsections, so numbering is section.index), computed mechanically by the lane's scanner.\par\medskip
\noindent\textit{Editorial note.} The counted proof contains no \textbackslash{}ref to any other labelled item of the paper, so this card needs no navigation stub at all; the only external inputs are the bibliography citations, which is why the publisher bibliography is attached.\par\medskip
\noindent\textit{Allowed standard prerequisites (editorial).}\begin{itemize}
\item Galois cohomology of a torus split by a finite Galois covering, via the Hochschild-Serre spectral sequence. Used inside the counted proof at lines 578-591. Produces the group G = Gal(V/U) and the exact sequence 0 -> H\textasciicircum{}1(G, T(V)) -> H\textasciicircum{}1(U, T) -> H\textasciicircum{}1(V, T\_V).
\item Finiteness of the ideal class group and Dirichlet's unit theorem over a global field. Used at lines 582 and 593-595. Makes H\textasciicircum{}0(V, G) of finite type and H\textasciicircum{}1(V, G) of finite type (finite when U is not X).
\item Artin-Mazur-Milne duality for fppf cohomology. Cited at lines 622-623. Transfers the finiteness of H\textasciicircum{}1(U, N) to H\textasciicircum{}2\_c(U, N\textasciicircum{}D).
\end{itemize}\medskip
\begin{lem} \label{fini1}
\begin{enumeratea}
\item\label{fini1a}
Let $\cT$ be a torus over $U$. Then
$H^0(U,\cT)$ and $H^1(U,\mathcal T)$
are of finite type. If~$U \neq X$, then
$H^1(U,\cT)$ is finite.

\item\label{fini1b}
Let $\cN$ be a finite group scheme of multiplicative type over
$U$. Then $H^1(U,\cN)$ and $H^2_\rc(U,\cN^D)$ are finite.
\end{enumeratea}
\end{lem}

\skpt\begin{proof}
\begin{enumeratea}
\item
Let $V$ be an étale and finite connected
Galois covering of $U$ such that $\cT_V:=\cT \times_U V$
is isomorphic to $\G^r$ for some non negative integer $r$.
The group $\cT(V) \simeq
H^0(V,\G)^r$ is of finite type by Dirichlet's theorem on units. Therefore
$H^0(U,\cT) \subset H^0(V,\cT)=\cT(V)$
is also of finite type.

Set $G=\gal(V/U)$. Then the Hochschild-Serre spectral sequence provides an
exact sequence
\[
0 \to H^1(G,\cT(V)) \to H^1(U,\cT)
\to H^1(V,\cT_V).
\]
Since $\cT_V \cong \G^r$, the group
$H^1(V,\cT_V) \cong (\pic V)^r$ is of finite type
(\resp finite if $U \neq X$) by finiteness of the ideal class group
of a global field.
As $\cT(V)$ is of finite type and $G$ is finite, the group
$H^1(G,\cT(V))$ is finite by \cite[Chap.\,VIII, Cor.\,2]{col}.
Thus $H^1(U,\cT)$ is of
finite type (\resp finite if $U \neq X$) as well.

\item
We can assume (by (\cite[Lem.\,III.8.9]{MilADT}
and \cite[Cor.\,4.9]{DemH}) that $U \neq X$.
Since every finitely generated Galois module is a quotient of a
torsion-free and finitely generated Galois module, the
assumption that $\cN$ is of multiplicative type implies
(by~\cite[X, Prop.\,1.1]{sga3})
that there is an exact sequence of $U$-group schemes
\[
0 \to \cN \to \cT_1 \to \cT_2 \to 0,
\]
where $\cT_1$ and $\cT_2$ are $U$-tori. Therefore, there is
an exact sequence of abelian groups
\[
H^0(U,\cT_2) \to H^1(U, \cN) \to H^1(U, \cT_1).
\]
By \eqref{fini1a}, we know that $H^1(U, \cT_1)$ is finite. Let $n$ be the order
of $\cN$; then the map $H^0(U,\cT_2) \to H^1(U, \cN)$
factorizes through a map $H^0(U,\cT_2)/n \to H^1(U, \cN)$.
But $H^0(U,\cT_2)/n$ is finite because $H^0(U,\cT_2)$
is of finite type by \eqref{fini1a}. Finally $H^1(U, \cN)$ is finite. The
finiteness of $H^2_\rc(U,N^D)$ follows by Artin-Mazur-Milne
duality (\cite[Th.\,1.1]{DemH}).\qedhere
\end{enumeratea}
\end{proof}

\begin{thebibliography}{99}
\bibitem{BD} M. Borovoi \& C. Demarche - “Manin obstruction to strong approximation for homogeneous spaces” , Comment. Math. Helv. 88 (2013) no. 1, p. 1-54
\bibitem{Bor} M. Borovoi - “The Brauer-Manin obstructions for homogeneous spaces with connected or abelian stabilizer” , J. reine angew. Math. 473 (1996), p. 181-194
\bibitem{bkiac} N. Bourbaki - Éléments de mathématique. Algèbre commutative. Chapitre 10 , Springer-Verlag, Berlin, 2007
\bibitem{Ces} K. Česnavičius - “Topology on cohomology of local fields” , Forum Math. Sigma 3 (2015), article ID e16, 55 pages
\bibitem{cesnalms} K. Česnavičius - “$p$-Selmer growth in extensions of degree $p$” , J. London Math. Soc. (2) 95 (2017) no. 3, p. 833-852
\bibitem{demlms} C. Demarche - “Le défaut d’approximation forte dans les groupes linéaires connexes” , Proc. London Math. Soc. (3) 102 (2011) no. 3, p. 563-597
\bibitem{demphd} C. Demarche - “Suites de Poitou-Tate pour les complexes de tores à deux termes” , Internat. Math. Res. Notices (2011) no. 1, p. 135-174, short version of “Théorèmes de dualité pour les complexes de tores”, available at https://webusers.imj-prg.fr/\textasciitilde{}cyril.demarche/
\bibitem{DemH} C. Demarche \& D. Harari - “Artin-Mazur-Milne duality for fppf cohomology” , Algebra Number Theory 13 (2019) no. 10, p. 2323-2357
\bibitem{friedsus} E. M. Friedlander \& A. Suslin - “The spectral sequence relating algebraic $K$-theory to motivic cohomology” , Ann. Sci. École Norm. Sup. (4) 35 (2002) no. 6, p. 773-875
\bibitem{gonz1} C. D. González-Avilés - “Arithmetic duality theorems for 1-motives over function fields” , J. reine angew. Math. 632 (2009), p. 203-231
\bibitem{gonz2} C. D. González-Avilés \& K.-S. Tan - “The generalized Cassels-Tate dual exact sequence for 1-motives” , Math. Res. Lett. 16 (2009) no. 5, p. 827-839
\bibitem{God} R. Godement - Topologie algébrique et théorie des faisceaux , Hermann, Paris, 1973
\bibitem{hss} D. Harari, C. Scheiderer \& T. Szamuely - “Weak approximation for tori over $p$-adic function fields” , Internat. Math. Res. Notices (2015) no. 10, p. 2751-2783
\bibitem{lang} S. Lang - “Algebraic groups over finite fields” , Amer. J. Math. 78 (1956), p. 555-563
\bibitem{MilEC} J. S. Milne - Étale cohomology , Princeton Math. Series , vol. 33 , Princeton University Press, Princeton, N.J., 1980
\bibitem{MilADT} J. S. Milne - Arithmetic duality theorems , BookSurge, LLC, Charleston, SC, 2006
\bibitem{nsw} J. Neukirch, A. Schmidt \& K. Wingberg - Cohomology of number fields , Grundlehren Math. Wiss. , vol. 323 , Springer-Verlag, Berlin, 2008
\bibitem{Ros} Z. Rosengarten - “Tamagawa numbers and other invariants of pseudo-reductive groups over global function fields” , 2018, to appear in Algebra Number Theory
\bibitem{San} J.-J. Sansuc - “Groupe de Brauer et arithmétique des groupes algébriques linéaires sur un corps de nombres” , J. reine angew. Math. 327 (1981), p. 12-80
\bibitem{col} J.-P. Serre - Corps locaux , Publications de l’Université de Nancago , vol. VIII , Hermann, Paris, 1968
\bibitem{SGA4} M. Artin, A. Grothendieck \& J.-L. Verdier - Théorie des topos et cohomologie étale des schémas. Tome 3 , Lect. Notes in Math. , vol. 269, 270 \& 305 , Springer-Verlag, Berlin-New York, 1973, Séminaire de Géométrie Algébrique du Bois-Marie 1963–1964 (SGA 4)
\bibitem{sga3} M. Demazure \& A. Grothendieck - Schémas en groupes (SGA 3). Tome I. Propriétés générales des schémas en groupes , Documents mathématiques , vol. 7 , Société Mathématique de France, Paris, 2011, article ID e16, Revised and annotated edition of the 1970 original
\bibitem{SP} Stacks Project Authors - “The Stacks Project” , 2019, https://stacks.math.columbia.edu
\end{thebibliography}
\end{document}
